\section*{5.3.2 Lifting (author source pp.27--30)}
Given Lemma 5.3, we are going to prove the following lifting for the certificate complexity.
\begin{originalresult}{Theorem 5.7}
For every function $f:\bits^n\to\bits$, we have $C(f)=O(C_\oplus(f\comp\OR)^3)$.
\end{originalresult}
Throughout the proof of the theorem, we will think of $f$ as having inputs $x_1,x_2,\ldots,x_n$, and of block composition $f\comp\OR$ as having inputs $z_{1,1},z_{1,2},z_{2,1},z_{2,2},\ldots,z_{n,1},z_{n,2}$, such that $x_i$ corresponds to $z_{i,1}\vee z_{i,2}$. The twin of $z_{i,1}$ is $z_{i,2}$, and vice versa. We start the proof with the following lemma, another technical tool in this section. Together with Lemma 5.3, the proof of the lifting theorem is essentially a combination of these two tricks.
\begin{originalresult}{Lemma 5.8}
Suppose that $C$ is a parity certificate for $f\comp\OR$ of size $k$ whose minimum weight is at least $k$. Then there are $k$ indices $i_1,\ldots,i_k$ such that for every input $x$ with $x_{i_1}=\cdots=x_{i_k}=1$, we can find an input $z$ satisfying $C$ such that $\OR^n(z)=x$. In particular, there is an input $z$ satisfying $C$ such that $\OR^n(z)=(1,\ldots,1)$.
\end{originalresult}
\begin{proof}
Let $V+b$ be the affine subspace corresponding to $C$, where $V$ is a linear subspace (the row space of the linear system $Ax=b$ in Definition 5.2). We pick a basis $\ell_1,\ldots,\ell_k$ for $V$ and indices $(i_1,j_1),\ldots,(i_k,j_k)\in[n]\times[2]$ such that:
\begin{itemize}
\item $z_{i_t,j_t}$ is in the support of $\ell_t$.
\item $z_{i_s,1},z_{i_s,2}$ are not in the support of $\ell_t$ for $s<t$.
\end{itemize}
We build the basis by successively adding vectors to it, similar to the Gaussian elimination. At step $t$, we have already chosen basis vectors $\ell_1,\ldots,\ell_{t-1}$ and the corresponding indices. Consider all the vectors in $V$ whose support does not contain any of the variables $z_{i_1,j_1},\ldots,z_{i_{t-1},j_{t-1}}$. Since $V$ has minimum weight $k>t-1$, we can find a vector $\ell_t$ whose support contains a variable $z_{i_t,j_t}$ with $i_t\notin\{i_1,\ldots,i_{t-1}\}$ and add it to the basis.
Now let $x$ be an input satisfying $x_{i_1}=\cdots=x_{i_k}=1$. We set the twins of $z_{i_1,j_1},\ldots,z_{i_t,j_t}$ to 1, and set $z_{i,1},z_{i,2}$ for $i\ne i_1,\ldots,i_t$ in an arbitrary way subject to $z_{i,1}\vee z_{i,2}=x_i$. We then set $z_{i_t,j_t},\ldots,z_{i_1,j_1}$ (in this order) in such a way that the resulting input satisfies $C$.
\end{proof}
The next lemma that we need for the proof of the lifting result gives an upper bound on the minimum weight of a parity certificate in terms of the sizes of certificates for 0 and 1.
\begin{originalresult}{Lemma 5.9}
Suppose that $f'\comp\OR$ has a parity certificate $C'$ of size $k'$ which certifies $b$, and a parity certificate $C''$ of size $k''$ which certifies $\bar b$. Then the minimum weight of $C'$ is at most $k'+2(k'')^2$.
\end{originalresult}
\begin{proof}
First, observe that both $k'>0$ and $k''>0$ since otherwise $f'\comp\OR$ is constant and cannot have certificates for both 0 and 1.
Consider the case $k''=1$. The certificate $C''$ consists of one constraint $\ell=\alpha$. By Lemma 5.8, since $C''$ has size 1, some input $z$ satisfying $C''$ satisfies $\OR^n(z)=(1,\ldots,1)$, and so $f'(1,\ldots,1)=\bar b$. Suppose for the sake of contradiction that the minimum weight of $C'$ is at least $k'+2(k'')^2+1=k'+3>k'$. According to Lemma 5.8, some input $z$ satisfying $C'$ satisfies also $\OR^n(z)=(1,\ldots,1)$, and so $f'(1,\ldots,1)=b$, which is impossible. So the minimum weight of $C'$ is at most $k'+2$.
We proceed by induction on $k''$. Suppose for the sake of contradiction that the minimum weight of $C'$ is at least $k'+2(k'')^2+1>k'$. According to Lemma 5.8, some input $z$ satisfying $C'$ satisfies $\OR^n(z)=(1,\ldots,1)$, and so $f'(1,\ldots,1)=b$. If the minimum weight of $C''$ was also at least $k''$, then the same argument would imply that $f'(1,\ldots,1)=\bar b$. Therefore, the minimum weight of $C'$ is less than $k''$. Consider some constraint $\ell=\alpha$ of $C''$ of weight at most $k''$. Lets choose some partial assignment satisfying $\ell=\alpha$, and complete it to a blockwise partial assignment $\tau$ in an arbitrary way. Thus $\tau$ assigns at most $2k''$ variables. By Lemma 5.4, the assignment $\tau$ cannot contradict $C'$ since the minimum weight of $C'$ is greater than $k'+2(k'')^2>2k''$. Therefore $(f'\comp\OR)|_\tau$ still has a $b$-certificate $C'|_\tau$ of size $k'$ and a $\bar b$-certificate $C''|_\tau$ of size $k''-1$. By induction, we know that the minimum weight of $C'|_\tau$ is at most $k'+2(k''-1)^2$, and so the minimum weight of $C'$ is at most $k'+2(k''-1)^2+2k''<k'+2(k'')^2$.
\end{proof}
The following lemma applies the `minimal weight lemma' to the current setting.
\begin{originalresult}{Lemma 5.10}
Suppose that $f\comp\OR$ has a parity certificate at some input $z$ of size at most $k$. Let $K$ be a parameter. Suppose that the certificate complexity of $f$ at $\OR^n(z)$ is greater than $k\cdot K$. Then we can find a blockwise partial assignment $\sigma$ to the inputs of $f\comp\OR$, consistent with $z$, such that for some $k'\le k$:
\begin{itemize}
\item $(f\comp\OR)|_\sigma$ has a parity certificate of size $k'$ at $z|_\sigma$ whose minimum weight is at least $K$.
\item $C(f|_{\OR^n(\sigma)},\OR^n(z)|_{\OR^n(\sigma)})>k'K$.
\end{itemize}
\end{originalresult}
\begin{proof}
We apply Lemma 5.3 for a function $f$ and a gadget $g=\OR$.
\end{proof}
Now we have all the ingredients for proving Theorem 5.7.
\begin{proof}[Proof of Theorem 5.7]
Let $C_\oplus(f\comp\OR)=k$, and let $K$ be a parameter to be determined later. Suppose that $C(f)>kK$. Apply Lemma 5.10 to get a non-constant function $f'$, obtained by restricting $f$, such that $f'\comp\OR$ has a parity certificate of size $k'$ whose minimum weight is at least $K$ (the function $f'$ is non-constant since its certificate complexity is positive). Since $C_\oplus(f'\comp\OR)\le C_\oplus(f\comp\OR)$ and $f'$ is non-constant, Lemma 5.9 (with $k''\le k$) shows that $K\le k'+2(k'')^2\le k+2k^2$. By choosing $K=k+2k^2+1$ we reach a contradiction. Therefore $C(f)\le kK=O(k^3)$.
\end{proof}
